\documentclass[11pt,a4paper]{article}

% --- Encoding + fonts ---------------------------------------------
\usepackage[T1]{fontenc}
\usepackage[utf8]{inputenc}
\usepackage{lmodern}
\usepackage{microtype}
% Allow up to 3em of extra stretch on lines TeX can't otherwise
% break cleanly — kills sub-mm overfulls without affecting normal
% paragraphs.
\setlength{\emergencystretch}{3em}

% --- Math ---------------------------------------------------------
\usepackage{amsmath, amssymb, amsthm, amsfonts}

% --- Layout -------------------------------------------------------
\usepackage[a4paper,margin=2.5cm]{geometry}
\usepackage{booktabs}
\usepackage{multirow}
\usepackage{array}
\usepackage[hidelinks]{hyperref}
\usepackage[safe]{tipa}

% --- Theorem environments ----------------------------------------
\theoremstyle{plain}
\newtheorem{theorem}{Theorem}[section]
\newtheorem{proposition}[theorem]{Proposition}
\newtheorem{lemma}[theorem]{Lemma}
\newtheorem{corollary}[theorem]{Corollary}

\theoremstyle{definition}
\newtheorem{definition}[theorem]{Definition}
\newtheorem{assumption}[theorem]{Assumption}
\newtheorem{example}[theorem]{Example}

\theoremstyle{remark}
\newtheorem{remark}[theorem]{Remark}

% --- Notation -----------------------------------------------------
\newcommand{\E}{\mathbb{E}}
\newcommand{\PP}{\mathbb{P}}
\newcommand{\R}{\mathbb{R}}
\newcommand{\N}{\mathbb{N}}
\newcommand{\Z}{\mathbb{Z}}
\newcommand{\one}{\mathbf{1}}
\newcommand{\indep}{\perp\!\!\!\perp}
\newcommand{\eqdef}{\mathrel{\stackrel{\mathrm{def}}{=}}}
\newcommand{\T}{\mathsf{T}}
\DeclareMathOperator*{\argmin}{arg\,min}
\DeclareMathOperator*{\argmax}{arg\,max}
\DeclareMathOperator{\Var}{Var}
\DeclareMathOperator{\Cov}{Cov}
\DeclareMathOperator{\Tr}{Tr}
\DeclareMathOperator{\IRR}{IRR}
\DeclareMathOperator{\ATE}{ATE}
\DeclareMathOperator{\ATT}{ATT}
\DeclareMathOperator{\ATC}{ATC}

\title{The MRM Framework: A Multi-Source Statistical Foundation
       for Canadian Carceral, Police, and Oversight Data,
       Implemented as MRM modules in MORIE}

\author{Vansh Singh Ruhela\thanks{%
  Correspondence: \texttt{hadesllm@proton.me}.
  Code: MORIE package (Python + R), released open-source.}}

\date{2026-05-09}

\begin{document}

\maketitle

\begin{abstract}
This article introduces \emph{MRM} (McNamara-Ruhela-Medina), a
multi-source statistical framework that brings five distinct
Canadian carceral and oversight data streams together under a
coordinated set of estimators and primitives. The framework spans (i)
the open Ontario Tracking Information System (OTIS) provincial
restrictive-confinement microdata, (ii) the federal Structured
Intervention Unit Implementation Advisory Panel (SIU~IAP)
reports and Sprott--Doob--Iftene independent academic
replications of federal SIU operations, (iii) the Ontario
Special Investigations Unit (Ontario~SIU) police-oversight
reports ingested by an automated scraper that I developed, (iv)
the Toronto Police Service (TPS) open-data crime categories and
the Statistics Canada Crime Severity Index, and (v) federal
Corrections and Conditional Release Statistical Overview (CCRSO)
aggregate tables introduced into the record by Doob's
\emph{T-539-20} affidavit. MRM is implemented in the
open-source MORIE package as a set of MRM modules,
comprising a 10-estimator per-individual ensemble, a GEE
clustering grid for high-cardinality groupings, a Doob
$\chi^{2}$ family for aggregate contingency tables, and a
Mandela~Rules classifier that operates at both the federal and
provincial level. I establish MRM's notation, derive each
estimator in its MRM instantiation, document the empirical
replication of all five published $\chi^{2}$ statistics from
Sprott--Doob--Iftene (reproducing every value to within $0.01$),
and report a headline empirical finding: the provincial Ontario
Mandela ``torture''-classified proportion rises from $12.5\%$ in
fiscal 2023 to $20.6\%$ in fiscal 2025, exceeding the federal
SIU rate of $9.9\%$ that Sprott and Doob reported for fiscal
2019--2020.
\end{abstract}

\section{Introduction}
\label{sec:intro}

Canadian carceral and oversight data are produced by multiple
disjoint agencies under disjoint legal frameworks, each with its
own data dictionary, release cadence, and analytic conventions.
At the federal level, Correctional Service Canada (CSC) operates
Structured Intervention Units (SIUs) under the
\emph{Corrections and Conditional Release Act} (CCRA), with
oversight provided by Independent External Decision Makers
(IEDMs) and a federally-appointed Implementation Advisory Panel
(SIU~IAP). At the provincial level, Ontario operates
restrictive-confinement and segregation placements that are
recorded in OTIS, an open dataset maintained by the Ministry of
the Solicitor General. Police-oversight in Ontario is exercised
by the Special Investigations Unit (Ontario~SIU)~--- a separate
agency whose acronym collides with the federal SIU. Crime data
are released by the Toronto Police Service (TPS) under nine
open-data categories, and aggregate national flow data are
published annually by Public Safety Canada in the CCRSO. None of
these streams shares a primary key, none uses the same time
window, and most are released as PDFs or HTML rather than tabular
microdata.

The contribution of this paper is a multi-source statistical
framework, \emph{MRM} (McNamara-Ruhela-Medina; named for its
methodology lineage --- see \S\ref{sec:mrm}), that brings these
five data streams into a coordinated notation and a coordinated
set of estimators. MRM is implemented in
MORIE~\cite{Ruhela2026MORIE}\footnote{MORIE = Methods for
Observational Inference and Robust Analysis of Interventions in
Scientific Experimentation. The package is hosted at
\url{https://github.com/hadesllm/morie}.}, an open-source Python
and R package released alongside this paper, as a set of \emph{MRM
modules} --- the code-level instantiation of MRM's estimators on
the five data sources above.
The mathematical content of MRM is divided into:
\begin{itemize}
  \item An aggregate IRR (incidence-rate-ratio) family for
        Poisson- and negative-binomial-distributed counts with a
        log offset (\S\ref{sec:mrm-aggregate}).
  \item A per-individual 10-estimator causal ensemble covering
        the ATE, ATT, and ATC under unconfoundedness, double-machine
        learning, propensity-score matching, and AIPW with a
        SuperLearner (\S\ref{sec:mrm-perrow}).
  \item A generalized-estimating-equations (GEE) clustering grid
        with priority ordering across families and working
        correlations (\S\ref{sec:gee-grid}).
  \item A Doob $\chi^{2}$ family for testing independence in
        aggregate contingency tables, with the Cram\'{e}r's $V$
        effect size (\S\ref{sec:doob-chi2}).
  \item A Mandela Rules classifier that operates at both the
        federal level (with the full out-of-cell hour
        operationalization) and the provincial level (with a
        duration-only proxy that I make explicit)
        (\S\ref{sec:mandela}).
  \item A Crime Severity Index (CSI) integration that turns TPS
        nine-category counts into Statistics Canada CSI scores
        (\S\ref{sec:csi}).
  \item A Goffmanian institutional-churn analysis suite
        (cyclical, mortifying, embedding dimensions; intra-year
        transition matrices; GLMM-IRR) on OTIS individual-level
        files (\S\ref{sec:churn}).
  \item A spatial--temporal stack for TPS data including a
        general Hawkes self-exciting process, Ripley's $K$,
        Moran's $I$, and a Getis--Ord $G^{\ast}$ statistic
        (\S\ref{sec:tps-spatial}).
  \item An automated scraper for the Ontario~SIU
        director's-report corpus (\S\ref{sec:siu-ontario}).
\end{itemize}

I prove a small set of structural results
(Theorems~\ref{thm:rf-irr-asymptotics} and~\ref{thm:gee-priority})
and replicate every published $\chi^{2}$ statistic from
Sprott--Doob--Iftene's three independent reports
(\S\ref{sec:replication}). The headline empirical finding,
reported in \S\ref{sec:findings} and discussed in
\S\ref{sec:mandela}, is that the proportion of OTIS provincial
restrictive-confinement individuals whose aggregate stay crosses
the Mandela 15-day threshold has risen monotonically from
$12.5\%$ in fiscal 2023 to $20.6\%$ in fiscal 2025, exceeding the
federal SIU torture-classified proportion of $9.9\%$ that Sprott
and Doob~\cite{SprottDoob2021Torture} found for fiscal 2019--2020,
even after applying the most conservative duration-only proxy.

\section{Notation}
\label{sec:notation}

Throughout, $i \in \{1, \ldots, n\}$ indexes individuals,
$j \in \{1, \ldots, J\}$ indexes institutions, $r \in \mathcal{R}$
indexes Ontario administrative regions or CSC federal regions, and
$t \in \{1, \ldots, T\}$ indexes fiscal years. I let $T_{i}$ be a
binary or categorical treatment, $Y_{i} \in \N$ be a count
outcome, $X_{i} \in \R^{p}$ be a vector of covariates, and
$E_{i} > 0$ be a person-time exposure (always treated as a fixed
offset on the log scale). I write
$T_{i}(t)\in\{0,1\}$ for the potential outcome under treatment
$t$, and assume the standard SUTVA, consistency, and positivity
conditions throughout.

\begin{definition}[Datasets]
\label{def:datasets}
The MRM framework integrates five primary data sources:
\begin{enumerate}
  \item $\mathcal{D}_{\mathrm{OTIS}}$: the Ontario Tracking
        Information System aggregate placements file
        ($n \approx 7.7\times 10^{4}$ row-aggregations, three
        fiscal years).
  \item $\mathcal{D}_{\mathrm{SIU}\text{-}\mathrm{IAP}}$: the
        federal SIU Implementation Advisory Panel reports plus
        the four Sprott--Doob--Iftene CRIMSL/Schulich
        publications.
  \item $\mathcal{D}_{\mathrm{Ontario}\text{-}\mathrm{SIU}}$: the
        Ontario Special Investigations Unit director's-report
        corpus, ingested by the scraper described in
        \S\ref{sec:siu-ontario} into a 65-column canonical schema
        keyed on the SIU case number $c \in \mathrm{YY\text{-}XXX\text{-}NNN}$.
  \item $\mathcal{D}_{\mathrm{TPS}}$: the Toronto Police Service
        open-data crime categories,
        \begin{multline*}
          \mathcal{C} \;=\; \{\,
            \text{Assault},\;
            \text{Auto Theft},\;
            \text{Bicycle Theft},\;
            \text{Break and Enter},\\
            \text{Homicide},\;
            \text{Robbery},\;
            \text{Shooting},\;
            \text{Theft from Motor Vehicle},\;
            \text{Theft Over}\,\}.
        \end{multline*}
  \item $\mathcal{D}_{\mathrm{CCRSO}}$: the federal Corrections
        and Conditional Release Statistical Overview
        aggregate tables introduced into the public record via
        Doob's $T$-$539$-$20$ affidavit~\cite{Doob2020Affidavit}.
\end{enumerate}
\end{definition}

\begin{remark}[Two distinct SIUs]
\label{rmk:two-siu}
The acronym SIU is used for two unrelated agencies. The
\emph{federal Structured Intervention Unit} is a CSC
post-segregation regime created by Bill C-83 in 2019; analyses
that draw on it use the module \texttt{morie.siuiap}. The
\emph{Ontario Special Investigations Unit} is a provincial
police-oversight body that investigates police-related deaths,
serious injuries, and sexual-assault complaints; analyses that
draw on it use the module \texttt{morie.siu}. The two never
share data, and care must be taken in cross-references; I have
deliberately chosen module names that do not begin with a common
prefix in order to make this distinction visible at import time.
\end{remark}

\section{The aggregate IRR family}
\label{sec:mrm-aggregate}

Let $Y_{ij}$ denote the count outcome for unit $i$ in stratum
$j \in \{1, \ldots, J\}$ and let $E_{ij}$ denote the corresponding
exposure (person-years, placement-days, or analogous person-time
denominator). MRM's aggregate component fits the log-linear
count model
\begin{equation}
  \log \E[Y_{ij} \mid X_{ij}] \;=\; \log E_{ij}
   \;+\; \alpha_{j} \;+\; \beta^{\T} X_{ij},
  \label{eq:rf-aggregate-model}
\end{equation}
under either a Poisson or negative-binomial distribution for
$Y_{ij}$, with $\alpha_{j}$ a stratum fixed effect and $\beta$ the
regression coefficients of interest. The treatment effect of a
binary $T \in \{0,1\}$ is summarized by the
\emph{incidence-rate ratio}
\begin{equation}
  \IRR \;\eqdef\;
   \frac{\E[Y \mid T=1, X]/E_{T=1}}
        {\E[Y \mid T=0, X]/E_{T=0}}
   \;=\; \exp(\beta_{T}).
  \label{eq:irr}
\end{equation}

\begin{theorem}[Asymptotics of the aggregate IRR]
\label{thm:rf-irr-asymptotics}
Under the standard Poisson regression regularity conditions and
the negative-binomial extension of~\eqref{eq:rf-aggregate-model},
the maximum-likelihood estimator $\hat{\beta}_{T}$ is consistent
and asymptotically normal:
\begin{equation*}
  \sqrt{n}\,(\hat{\beta}_{T} - \beta_{T})
   \;\xrightarrow{d}\; \mathcal{N}\!\left(0,\, \mathbf{I}^{-1}_{TT}\right),
\end{equation*}
where $\mathbf{I}^{-1}_{TT}$ is the $(T,T)$ block of the inverse
Fisher information matrix. The Wald confidence interval for the
incidence-rate ratio at level $1 - \alpha$ is
\begin{equation}
  \left[\,
  \exp\!\left(\hat{\beta}_{T} - z_{1 - \alpha/2}
              \cdot \widehat{\mathrm{SE}}(\hat{\beta}_{T})\right),
  \;
  \exp\!\left(\hat{\beta}_{T} + z_{1 - \alpha/2}
              \cdot \widehat{\mathrm{SE}}(\hat{\beta}_{T})\right)
  \,\right].
  \label{eq:irr-ci}
\end{equation}
\end{theorem}

When the stratum count $J$ is high (e.g., $J = 25$ Ontario
correctional institutions), the per-stratum maximum-likelihood
estimator of $\alpha_{j}$ has a high effective dimension and the
finite-sample distribution is poorly approximated by the
asymptotic Normal. I address this in \S\ref{sec:gee-grid} via a
generalized-estimating-equations (GEE) approach with cluster-robust
inference. In MORIE, this aggregate component is implemented as
the aggregate MRM module.

\section{Per-individual estimators: the MRM (10 estimators)}
\label{sec:mrm-perrow}

Where individual-level microdata are available
($\mathcal{D}_{\mathrm{OTIS}}$ files \texttt{a01}, \texttt{b01},
\texttt{b02}), MRM targets the average treatment effect (ATE),
the effect on the treated (ATT), and the effect on the controls
(ATC):
\begin{align}
  \tau_{\mathrm{ATE}} &\eqdef
    \E\bigl[Y(1) - Y(0)\bigr],
  \label{eq:ate} \\
  \tau_{\mathrm{ATT}} &\eqdef
    \E\bigl[Y(1) - Y(0) \;\big|\; T = 1\bigr],
  \label{eq:att} \\
  \tau_{\mathrm{ATC}} &\eqdef
    \E\bigl[Y(1) - Y(0) \;\big|\; T = 0\bigr].
  \label{eq:atc}
\end{align}
I instantiate each via ten complementary estimators, collectively
implemented in MORIE as the \emph{Per-individual Ruhela
Formulations} (per-row MRM).

\subsection{Inverse-probability-weighting estimators}

Let $e(x) \eqdef \PP(T = 1 \mid X = x)$ denote the propensity
score, and $\hat{e}(x)$ a parametric or machine-learning
estimate. The H\'ajek-stabilized inverse-probability-weighting
ATE is
\begin{equation}
  \widehat{\tau}^{\mathrm{IPW}}_{\mathrm{ATE}}
  \;=\;
  \frac{\sum_{i : T_i = 1} Y_i / \hat{e}(X_i)}
       {\sum_{i : T_i = 1} 1 / \hat{e}(X_i)}
  \;-\;
  \frac{\sum_{i : T_i = 0} Y_i / (1 - \hat{e}(X_i))}
       {\sum_{i : T_i = 0} 1 / (1 - \hat{e}(X_i))}.
  \label{eq:ipw-ate}
\end{equation}
The corresponding ATT and ATC estimators reweight each side by
$\hat{e}(X_i)$ or $1 - \hat{e}(X_i)$ respectively.

\subsection{Augmented inverse-probability-weighting (AIPW)}

Let $\mu_{t}(x) \eqdef \E[Y \mid T=t, X=x]$ and let $\hat{\mu}_{t}$
be a regression estimate. The AIPW (doubly robust) estimator is
\begin{align}
  \widehat{\tau}^{\mathrm{AIPW}}_{\mathrm{ATE}}
  &= \frac{1}{n} \sum_{i=1}^{n}
   \biggl[ \hat{\mu}_{1}(X_i) - \hat{\mu}_{0}(X_i)
   \nonumber \\
  &\quad
   + \frac{T_i\,(Y_i - \hat{\mu}_{1}(X_i))}{\hat{e}(X_i)}
   - \frac{(1-T_i)(Y_i - \hat{\mu}_{0}(X_i))}{1 - \hat{e}(X_i)}
   \biggr].
  \label{eq:aipw-ate}
\end{align}
$\widehat{\tau}^{\mathrm{AIPW}}$ is consistent if either the
propensity model or the outcome model is correctly specified, and
attains the semiparametric efficiency bound when both are
correctly specified.

\subsection{Double-machine-learning interactive regression model (IRM-DML)}

Following Chernozhukov et al.~\cite{Chernozhukov2018DoubleML},
let $\hat{\mu}_{0}$ and $\hat{\mu}_{1}$ be cross-fitted estimates
of the conditional outcome and let $\hat{e}$ be a cross-fitted
propensity. The IRM-DML estimator is the empirical mean of the
score
\begin{equation}
  \psi(W; \tau, \mu_{0}, \mu_{1}, e)
  \;=\; \mu_{1}(X) - \mu_{0}(X)
   + \frac{T \cdot (Y - \mu_{1}(X))}{e(X)}
   - \frac{(1-T)(Y - \mu_{0}(X))}{1 - e(X)} - \tau,
  \label{eq:irm-score}
\end{equation}
which is Neyman-orthogonal in $(\mu_{0}, \mu_{1}, e)$. Cross-fitting
breaks the bias arising from over-fitting in the nuisance
estimators, and root-$n$ inference for $\tau$ is recovered even
when each nuisance is fit by an arbitrary learner converging at
the slower rate $n^{-1/4}$.

\subsection{Propensity-score matching and subclassification}

Let $d(x_{i}, x_{j})$ be a distance on covariates (typically the
logit of the propensity). The 1:1 nearest-neighbour matched ATE
estimator is
\begin{equation}
  \widehat{\tau}^{\mathrm{PSM}}_{\mathrm{ATE}} \;=\;
   \frac{1}{n} \sum_{i=1}^{n}
   \bigl[Y_{i} - Y_{m(i)}\bigr]
   \cdot (2T_{i} - 1),
  \label{eq:psm}
\end{equation}
where $m(i) \in \argmin_{j : T_{j} = 1 - T_{i}} d(X_{i}, X_{j})$.
The subclassification variant partitions the sample into $K$
strata of $\hat{e}$ and computes a stratum-weighted average of
within-stratum mean differences. I follow Rosenbaum and Rubin's
$K = 5$ default.

\subsection{Partial-linear and SuperLearner variants}

The partial-linear-regression (PLR) DML estimator drops the
interactive interaction term and is convenient when treatment
effects are believed homogeneous in $X$. The
AIPW-SuperLearner variant replaces the parametric $\hat{\mu}_{t}$
and $\hat{e}$ with the convex-combination
SuperLearner~\cite{vanderLaan2007SuperLearner}, which
asymptotically achieves the oracle risk of the best base learner.

\subsection{Canonical and dual MRM formulations on \texttt{a01},
            \texttt{b01}, and \texttt{b02}}
\label{sec:dual-rf}

In the MORIE instantiation, the per-row MRM instantiates the
ten estimators above on three OTIS individual-level files:
\texttt{a01} (restrictive-confinement person-rows), \texttt{b01}
(segregation person-rows), and \texttt{b02} (segregation
person-rows with the additional Toronto-region indicator). For
each, a canonical pair of treatment-and-outcome variables is
combined with a \emph{dual} naive-arm sensitivity pair to test
the identification claim. The full dual (canonical $+$ naive-arm
$+$ all alt-$T$ variants $+$ subgroup stratifications) is
implemented on both \texttt{a01} and \texttt{b01}; \texttt{b02}
adds an alt-region variant on top of the canonical formulation.

\paragraph{canonical MRM.}
Let $T^{\mathrm{ac}}_{i} \eqdef \one\{\mathrm{alert\_state}_{i}
\in \mathcal{S}_{\mathrm{high}}\}$, where
$\mathcal{S}_{\mathrm{high}}$ is the set of high-severity
alert-state combinations (mental-health alert, suicide-risk
alert, or suicide-watch alert flagged for individual $i$). Let
$Y^{\mathrm{vm}}_{i}$ denote the count of restrictive-confinement
\emph{movements} for individual $i$ within the fiscal year. The
canonical MRM is
\begin{equation}
  \tau^{\mathrm{(can)}}
  \;\eqdef\;
   \E\!\bigl[\,Y^{\mathrm{vm}}(1) - Y^{\mathrm{vm}}(0)\;\big|\;
              X = x\,\bigr],
   \qquad T = T^{\mathrm{ac}},
  \label{eq:rf-canonical}
\end{equation}
estimated via every estimator in
\S\ref{sec:mrm-perrow}.

\paragraph{dual (naive-arm) MRM.}
The naive arm replaces the high-alert binary with the
\emph{any-flag} binary
$T^{\mathrm{any}}_{i} \eqdef \one\{\sum_{s} \mathrm{flag}_{is}
\ge 1\}$ and the count outcome with the binary
$Y^{\mathrm{vm,bin}}_{i} \eqdef \one\{Y^{\mathrm{vm}}_{i} \ge 1\}$,
\begin{equation}
  \tau^{\mathrm{(dual)}}
  \;\eqdef\;
   \E\!\bigl[\,Y^{\mathrm{vm,bin}}(1) - Y^{\mathrm{vm,bin}}(0)\;\big|\;
              X = x\,\bigr],
   \qquad T = T^{\mathrm{any}}.
  \label{eq:rf-dual}
\end{equation}
Identification of $\tau^{\mathrm{(can)}}$ relies on the high-alert
indicator's discrimination among the alert-state taxonomy;
$\tau^{\mathrm{(dual)}}$ relies on the cruder ``any-flag versus
no-flag'' distinction. A pattern of agreement between
$\tau^{\mathrm{(can)}}$ and $\tau^{\mathrm{(dual)}}$ (both
significant, same sign, similar magnitude) is taken as
sensitivity-supporting evidence; a pattern of disagreement is a
flag that the identification rests on the finer-grained alert-
state distinction and should be reported with care.

\paragraph{Alt-$T$ variants.}
For each of the canonical and dual RFs, the framework reports
three alternative-treatment specifications by replacing
$T^{\mathrm{ac}}$ in turn with
\begin{itemize}
  \item $T^{\mathrm{gender}}_{i} = \one\{\mathrm{gender}_{i} = \mathrm{Female}\}$,
  \item $T^{\mathrm{age50+}}_{i} = \one\{\mathrm{age\_category}_{i} \ge 50\}$,
  \item $T^{\mathrm{Toronto}}_{i} = \one\{\mathrm{region}_{i} = \mathrm{Toronto}\}$,
\end{itemize}
and additionally fits subgroup-stratified variants on
$\{\mathrm{Female}\}$ and $\{\mathrm{Male}\}$ separately to test
effect heterogeneity.

\subsection{Sensitivity analysis: the E-value}

I follow Vanderweele and Ding~\cite{VanderWeele2017Evalue} and
report, for every per-individual estimate, the E-value
\begin{equation}
  E\text{-value}\;=\; \widehat{\IRR} + \sqrt{\widehat{\IRR}\,
   (\widehat{\IRR} - 1)},
  \label{eq:evalue}
\end{equation}
which is the smallest unmeasured-confounder strength (on the IRR
scale, jointly with both $T$ and $Y$) that would be required to
explain the observed estimate away to the null.

\section{GEE clustering grid}
\label{sec:gee-grid}

When the stratification implied by~\eqref{eq:rf-aggregate-model}
has high cardinality (e.g., $J = 25$ Ontario institutions or
$R = 5$ federal CSC regions), classical Poisson and
negative-binomial maximum-likelihood inference is fragile. MRM
instead estimates a marginal model
\begin{equation}
  g\bigl(\E[Y_{ij} \mid X_{ij}]\bigr) \;=\;
   \log E_{ij} \;+\; \beta^{\T} X_{ij},
  \label{eq:rf-gee-marginal}
\end{equation}
where $g$ is the log link, and supplies the within-cluster
correlation through a working-correlation matrix
$\mathbf{W}(\rho)$. Inference is then conducted via cluster-robust
sandwich variance.

\begin{theorem}[GEE priority ordering]
\label{thm:gee-priority}
Among the four combinations of family
$F \in \{\text{Poisson}, \text{NB}\}$ and working correlation
$\mathbf{W} \in \{\text{Exch}, \text{Indep}\}$, MRM selects the
converged fit with priority order
\begin{equation*}
  (\text{NB}, \text{Exch}) \succ (\text{Pois}, \text{Exch})
   \succ (\text{NB}, \text{Indep}) \succ (\text{Pois}, \text{Indep}),
\end{equation*}
breaking ties by minimum quasi-information criterion (QIC).
The chosen estimator is the most flexible converged option,
balancing overdispersion accommodation against working-correlation
sophistication.
\end{theorem}

In practice I report all four rows in the framework's output
table for sensitivity analysis; under the priority ordering, the
top-row estimate is taken as the canonical MRM.

\section{The MRM framework}
\label{sec:mrm}

MRM is the master abstraction this paper introduces: a multi-source
statistical framework for analysing Canadian carceral, police, and
oversight data, comprising the aggregate IRR family of
\S\ref{sec:mrm-aggregate}, the per-individual 10-estimator
ensemble of \S\ref{sec:mrm-perrow}, the GEE clustering grid of
\S\ref{sec:gee-grid}, the Doob $\chi^{2}$ family of
\S\ref{sec:doob-chi2}, and the Mandela--Rules classifier of
\S\ref{sec:mandela}, together with the data-pipeline and
descriptive components of
\S\S\ref{sec:csi}--\ref{sec:overlay}. The acronym names the
methodology lineage as catalyst--principal--supervisor:
\begin{itemize}
  \item \textbf{M --- McNamara (catalyst).} Glenn McNamara
        contributed the methodological foundations that catalysed
        this work: distribution theory, applied statistics for
        administrative data, and the statistical-judgment intuition
        drawn from a 35-year career with the Ontario Government.
        His mentorship grounded much of the framework's
        empirical-statistical practice.
  \item \textbf{R --- Ruhela (principal).} The framework author:
        the MRM notation, the per-row MRM (the 10 estimators of
        \S\ref{sec:mrm-perrow}), the cross-jurisdiction Mandela
        classifier (\S\ref{sec:mandela}), the Ontario~SIU
        automation pipeline (\S\ref{sec:siu-ontario}), the TPS--CSI
        integration (\S\ref{sec:csi}), and the MRM modules
        implementation in MORIE that carries this paper's
        analyses end-to-end.
  \item \textbf{M --- Medina (supervisor \& reviewer).} Prof.\
        Angela Zorro Medina, of the Centre for Criminology and
        Sociolegal Studies, is the author's supervisor and
        methodological instructor, served as the domain-expert
        reviewer of the preliminary methodological approach, and is
        a knowledge user of this framework. Her own work on
        anti-gang
        legislation~\cite{ZorroMedina2023JMP} establishes the
        methodological lineage MRM follows on the U.S./Canada
        bridge: a staggered two-way-fixed-effects identification
        strategy with formal leads-and-lags Granger-causality
        diagnostics for the parallel-trends
        assumption~\cite{Athey2022Staggered, CallawaySantanna2021,
        GoodmanBacon2021DiD}; an explicitly multi-source
        data-integration pattern (her five U.S.\ sources --- FBI
        SRS/NIBRS, BJS BJSPS, FBI UCR, BLS, state annotated
        criminal codes --- are structurally analogous to MRM's
        five Canadian sources of \S\ref{sec:notation}); the
        deterrence~/ routine-activities~/ certainty mechanism
        categorisation that grounds individual estimands in
        sociolegal theory; and the inequality-effects-of-criminal
        -law-expansions framing (\S2.3 of \cite{ZorroMedina2023JMP})
        that connects empirical estimands to racial and social
        inequality in carceral outcomes. Her substantive review of
        the OTIS Major Research Paper (Paper 204) --- insisting on
        quantitatively-grounded sociolegal mechanism, on a
        negative-binomial mixed model with regional-cluster random
        intercepts as the principal aggregate-level estimator, and
        on a disciplined empirical structure --- directly shapes
        the methodology reported here.
\end{itemize}

MRM is the framework; the \emph{MRM modules} are its MORIE
instantiation --- the Python and R code that implements every MRM
estimator on the five data sources of \S\ref{sec:notation}. The
MRM acronym is a methodology-attribution chain, not a
co-authorship statement: the framework, code, and findings
reported in this paper are the work of the framework author and
should be cited as~\cite{Ruhela2026MRM}.

\section{Doob $\chi^{2}$ family for aggregate contingency tables}
\label{sec:doob-chi2}

For an $r \times c$ contingency table with cell counts
$O_{ij}$ and row/column totals $R_{i}$ and $C_{j}$, the Pearson
$\chi^{2}$ statistic is
\begin{equation}
  \chi^{2}
  \;=\;
  \sum_{i=1}^{r} \sum_{j=1}^{c}
  \frac{\bigl(O_{ij} - E_{ij}\bigr)^{2}}{E_{ij}},
  \quad
  E_{ij} = R_{i} C_{j} / N,
  \label{eq:chi2}
\end{equation}
where $N = \sum_{ij} O_{ij}$. Under the null hypothesis of
independence, $\chi^{2}$ is asymptotically chi-square distributed
with $(r-1)(c-1)$ degrees of freedom. The accompanying effect-size
measure is Cram\'er's $V$:
\begin{equation}
  V \;=\; \sqrt{\frac{\chi^{2}/N}{\min(r-1, c-1)}}.
  \label{eq:cramers-v}
\end{equation}
MRM's Doob $\chi^{2}$ family applies \eqref{eq:chi2} and
\eqref{eq:cramers-v} to every meaningful 2-way slice of the OTIS
aggregate tables (region $\times$ gender, region $\times$ year,
institution $\times$ alert state, $\ldots$), together with
year-over-year trend tests on death counts.

\section{Mandela Rules classifier}
\label{sec:mandela}

The Mandela Rules~\cite{UN2015Mandela} (UN Standard Minimum Rules
for the Treatment of Prisoners) define solitary confinement and
torture under Rules 43 and 44. Rule 44 defines solitary
confinement as confinement for $22+$ hours per day without
meaningful human contact; Rule 43 prohibits the prolonged version
of this practice (defined as exceeding $15$ consecutive days) as
torture or other cruel, inhuman, or degrading treatment.

\subsection{Federal operationalization}

Sprott and Doob~\cite{SprottDoob2021Torture} operationalized
Rules 43 and 44 jointly on federal SIU person-stays:
\begin{equation}
  \mathrm{Class}(d, h, m)
  =
  \begin{cases}
    \text{Torture}    & \text{if } d \ge 16 \text{ and }
                              h \le 2 \text{ and } m = 100,\\
    \text{Solitary}   & \text{if } d \le 15 \text{ and }
                              h \le 2 \text{ and } m = 100,\\
    \text{All other}  & \text{otherwise},
  \end{cases}
  \label{eq:mandela-class}
\end{equation}
where $d$ is days in SIU, $h$ is the average hours-out-of-cell
per day, and $m \in [0,100]$ is the percentage of days on which
the inmate did not receive the legally-required four hours
out-of-cell. Across $N = 1{,}960$ federal SIU person-stays
(November 2019 -- September 2020), the classifier yields
$28.4\%$ \emph{Solitary}, $9.9\%$ \emph{Torture}, and $61.7\%$
\emph{All other}~\cite{SprottDoob2021Torture}.

\subsection{Provincial operationalization}

OTIS does not expose hours-out-of-cell, so MRM applies a
duration-only proxy at the provincial level:
\begin{equation}
  \mathrm{Class}^{\ast}(d) =
  \begin{cases}
    \text{Torture}^{\ast}  & \text{if } d \ge 16,\\
    \text{Solitary}^{\ast} & \text{if } 1 \le d \le 15.
  \end{cases}
  \label{eq:mandela-prov}
\end{equation}
The proxy is an upper bound on the true Mandela classification:
some placements with $d \ge 16$ may, in fact, provide more than
two hours out-of-cell on average, in which case the formal
Mandela definition is not met. I am explicit about this in
every analyzer's docstring, and I report both the
\emph{Segregation} and the broader \emph{Restrictive Confinement}
views of the OTIS \texttt{c11} aggregate.

\subsection{Cross-jurisdiction comparison}

Table~\ref{tab:mandela-cross} reports the Mandela
classification for the federal SIU regime~\cite{SprottDoob2021Torture}
alongside the provincial OTIS Segregation and Restrictive
Confinement views, by fiscal year. The provincial Segregation
proportion classified as $\text{Torture}^{\ast}$ rises from
$12.5\%$ in fiscal 2023 to $20.6\%$ in fiscal 2025 --- a peak gap
of $+10.7$ percentage points relative to the federal $9.9\%$
that Sprott and Doob reported.

\begin{table}[ht]
  \caption{Cross-jurisdiction Mandela classification.
           Federal: person-stays, full operationalization
           (\eqref{eq:mandela-class}). Provincial: individuals,
           duration-only proxy (\eqref{eq:mandela-prov}).
           Source: \cite{SprottDoob2021Torture} and the OTIS
           \texttt{c11} aggregate.}
  \label{tab:mandela-cross}
  \centering
  \begin{tabular}{lllrrr}
    \toprule
    Source & Period & Unit & Solitary \% & Torture \% & $N$ \\
    \midrule
    Federal SIUs        & 2019--2020 & Person-stays & 28.4 & 9.9  & 1{,}960 \\
    Provincial Seg.\    & 2023       & Individuals & 87.5 & 12.5 & 12{,}647 \\
    Provincial Seg.\    & 2024       & Individuals & 83.5 & 16.5 & 10{,}881 \\
    Provincial Seg.\    & 2025       & Individuals & 79.4 & 20.6 &  9{,}608 \\
    Provincial RC       & 2023       & Individuals & 68.5 & 31.5 & 20{,}781 \\
    Provincial RC       & 2024       & Individuals & 64.0 & 36.0 & 19{,}641 \\
    Provincial RC       & 2025       & Individuals & 59.1 & 40.9 & 25{,}045 \\
    \bottomrule
  \end{tabular}
\end{table}

\subsection{Caveats}

The federal--provincial comparison in
Table~\ref{tab:mandela-cross} is directional rather than
apples-to-apples. The federal numbers count person-stays and use
the full out-of-cell-hour operationalization; the provincial
numbers count individuals (which integrate across multiple
placements within a fiscal year) and use the duration-only proxy.
Two consequences follow. First, the provincial denominators are
larger and therefore the provincial denominator is a different
unit from the federal. Second, the provincial numerator is an
\emph{upper bound} on the formal Mandela classification, since
some placements with $d \ge 16$ may have provided more than two
hours out-of-cell on average. I have nevertheless reported the
proxy alongside the federal figures because the magnitude of the
gap, even after these conservative adjustments, is substantial.

\section{Crime Severity Index integration}
\label{sec:csi}

The Statistics Canada Crime Severity
Index~\cite{Wallace2009CSI} weights each \emph{Criminal Code}
offence by a sentence-and-incarceration severity:
\begin{equation}
  w(o) \;=\; \mathrm{avg}\bigl(\text{sentence days} \mid \text{convicted of } o\bigr)
   \cdot \PP(\text{incarcerated} \mid \text{convicted of } o).
  \label{eq:csi-weight}
\end{equation}
MRM integrates the CSI weights for the nine TPS open-data
categories. Per-year and per-neighbourhood CSI scores are then
\begin{equation}
  \mathrm{CSI}_{r, t}
  \;=\;
  \frac{\displaystyle\sum_{o \in \mathcal{C}}
   w(o) \cdot N_{r, t}(o)}
       {\mathrm{Pop}_{r, t}}\, \cdot 10^{5},
  \label{eq:csi-rate}
\end{equation}
where $r$ is a neighbourhood (HOOD\_158 region in TPS), $t$ is a
calendar year, $N_{r,t}(o)$ is the count of incidents of category
$o$, and the multiplier $10^{5}$ standardizes the rate to a
per-100,000-residents scale.

\section{Stochastic physics of crime}
\label{sec:tps-physics}

MRM integrates a small set of stochastic-physics models that
have appeared in the urban-crime
literature~\cite{Bettencourt2007, ShortBrantingham2008}. They are
implemented in \texttt{morie.tps\_statphysics} and provide a
mechanistic interpretation that complements the descriptive
spatial--temporal statistics of \S\ref{sec:tps-spatial}.

\subsection{Reaction--diffusion of crime density}

Let $u(x, t)$ denote the smoothed density of incidents of a given
TPS category at location $x$ and time $t$. MRM fits the
reaction--diffusion process
\begin{equation}
  \frac{\partial u}{\partial t}
  \;=\; D \nabla^{2} u
   \;+\; \omega\, u\,(1 - u)
   \;-\; \eta\, p(x, t),
  \label{eq:reaction-diffusion}
\end{equation}
where $D$ is the diffusion coefficient, $\omega$ the intrinsic
growth rate, $p(x, t)$ a police-attention field, and $\eta$ the
attention coupling strength.
Equation~\eqref{eq:reaction-diffusion} is the
Short--Brantingham--D'Orsogna model adapted to the TPS
neighbourhood grid; I discretise on a $90 \times 60$ grid
covering the city of Toronto.

\subsection{L\'evy-flight scaling}

Empirical inter-event spatial displacement is modelled by a
heavy-tailed jump distribution. I fit the L\'evy tail-index
$\alpha$ in
\begin{equation}
  \PP(\lvert \Delta x \rvert > r) \;\propto\; r^{-\alpha},
  \quad r \ge r_{\min},
  \label{eq:levy-tail}
\end{equation}
via the Hill estimator and report the scaling regime
$1 < \alpha < 3$ (super-diffusive) versus $\alpha \ge 3$
(diffusive) per category.

\subsection{Urban-scaling exponent}

Across cities (or, in this setting, across Toronto neighbourhoods
of varying populations), the urban-scaling
hypothesis~\cite{Bettencourt2007} predicts
\begin{equation}
  Y \;\propto\; \mathrm{Pop}^{\beta},
  \label{eq:urban-scaling}
\end{equation}
where $Y$ is a per-place crime count, $\beta < 1$ implies
sublinear scaling (per-capita crime falls with size), $\beta = 1$
linear, and $\beta > 1$ superlinear. MRM fits $\beta$ by
ordinary least squares on $\log Y \sim \log \mathrm{Pop}$ with
an HC3 covariance estimator.

\subsection{Lotka--Volterra police--crime dynamics}

For paired annual time series $C_{t}$ (crimes) and $P_{t}$
(police-officer count) within a category, the framework fits
\begin{align}
  \frac{dC}{dt} &= a\,C - b\,C\,P, \\
  \frac{dP}{dt} &= -c\,P + d\,C\,P,
  \label{eq:lotka-volterra}
\end{align}
the standard predator--prey system. The resulting phase portrait
gives a structural complement to the descriptive year-over-year
regression of \S\ref{sec:overlay}.

\section{TPS spatial-temporal stack}
\label{sec:tps-spatial}

\subsection{Hawkes self-exciting process}

For an event time series $\{t_{1}, \ldots, t_{n}\} \subset [0, T]$
within a TPS category, MRM fits the marked Hawkes process with
conditional intensity
\begin{equation}
  \lambda(t \mid \mathcal{H}_{t})
  \;=\; \mu(t) \;+\;
  \sum_{i \,:\, t_{i} < t} \alpha\, \kappa(t - t_{i}; \theta),
  \label{eq:hawkes}
\end{equation}
where $\alpha > 0$ is the branching ratio (excitation amplitude),
$\mu(\cdot)$ is a parametric baseline, and
$\kappa(\cdot; \theta)$ is the excitation kernel.

\paragraph{Baselines.} MRM supports two baselines: a constant
$\mu(t) = \mu_{0}$ and a sinusoidal
$\mu(t) = \mu_{0}\,\bigl[1 + \mu_{1}\sin(2\pi t/365 + \phi)\bigr]$
that captures yearly seasonality.

\paragraph{Kernels.} Four excitation kernels are supported:
\begin{align}
  \kappa^{\mathrm{exp}}(u; \beta)
   &\;=\; \beta\, e^{-\beta u},
   & u &\ge 0,
   \label{eq:hawkes-exp} \\
  \kappa^{\mathrm{gam}}(u; \alpha, \beta)
   &\;=\; \frac{\beta^{\alpha}}{\Gamma(\alpha)}\, u^{\alpha - 1}\, e^{-\beta u},
   & u &\ge 0,
   \label{eq:hawkes-gam} \\
  \kappa^{\mathrm{wei}}(u; \alpha, \lambda)
   &\;=\; \frac{\alpha}{\lambda}\!\left(\frac{u}{\lambda}\right)^{\!\alpha-1}
       \!e^{-(u/\lambda)^{\alpha}},
   & u &\ge 0,
   \label{eq:hawkes-wei} \\
  \kappa^{\mathrm{lom}}(u; \alpha, c)
   &\;=\; \frac{(\alpha - 1)\,c^{\alpha - 1}}{(u + c)^{\alpha}},
   & u &\ge 0,\;\alpha > 1.
   \label{eq:hawkes-lom}
\end{align}
The exponential kernel~\eqref{eq:hawkes-exp} is the classical
Markovian Hawkes case; the Lomax (power-law)
kernel~\eqref{eq:hawkes-lom} produces a non-Markovian process
with long-memory tails relevant to repeat-victimization dynamics.
The Weibull and gamma kernels span the intermediate regime.

\paragraph{Estimation and selection.} Each
$(\mu, \alpha, \theta)$ tuple is fit by maximum likelihood on
the conditional-intensity log-likelihood, and the framework
selects among kernels by AIC/BIC and a time-rescaling residual
Q--Q diagnostic~\cite{Daley2003PointProcess}.  A self-contained
treatment of the likelihood, score, and information for the
non-stationary, non-Markovian case --- including the
Kwan--Chen--Dunsmuir asymptotic-ergodicity argument that secures
consistency and asymptotic normality of the MLE when the kernel
is non-exponential and the baseline is non-stationary --- and the
Toronto Police Service application of the same eight (kernel
$\times$ baseline) combinations is given in the companion
paper~\cite{Ruhela2026Hawkes}.

\subsection{Spatial cross-correlation}

The bivariate Moran's $I$ statistic between two TPS categories
$a, b \in \mathcal{C}$ on a neighbourhood graph with row-normalized
adjacency $\mathbf{W}$ is
\begin{equation}
  I_{a,b}
  \;=\;
  \frac{n \sum_{i,j} W_{ij}
        (z^{a}_{i})(z^{b}_{j})}
       {\sum_{i,j} W_{ij}
        \cdot \sqrt{\sum_i (z^a_i)^2}
        \cdot \sqrt{\sum_i (z^b_i)^2}},
  \label{eq:bivariate-morans}
\end{equation}
where $z^{a}_{i}$ and $z^{b}_{i}$ are within-category z-scores of
neighbourhood-level counts. Local Moran's $I_{i}$, Getis--Ord
$G^{\ast}_{i}$, and Ripley's $K(r)$ follow standard definitions
and are exposed in the framework.

\section{Goffmanian institutional churn}
\label{sec:churn}

MRM operationalizes Goffman's
\emph{total institutions}~\cite{Goffman1961Asylums} thesis as a
suite of formal statistical tests on OTIS individual-level files.
Goffman's argument is that carceral environments produce a
cyclical, mortifying, and embedding relationship between the
institution and the individual; I instantiate each of those
three dimensions as a measurable quantity. The implementation
lives in \texttt{morie.otis\_churn}.

\paragraph{Cyclical dimension.}
Repeat-placement concentration is summarised by the Gini
coefficient on the within-year per-individual placement count
$\{N_{i}\}_{i=1}^{n}$:
\begin{equation}
  G \;=\; \frac{\sum_{i=1}^{n}\sum_{j=1}^{n}
                |N_{i} - N_{j}|}
              {2 n^{2} \bar{N}},
  \quad
  \bar{N} = n^{-1}\sum_{i} N_{i}.
  \label{eq:churn-gini}
\end{equation}
MRM also fits a power-law tail
$\PP(N \ge k) \propto k^{-\gamma}$ via maximum likelihood and tests
its goodness of fit by a Kolmogorov--Smirnov statistic, after
\cite{ClausetShalizi2009Powerlaw}. Within-year region diversity ---
how many distinct administrative regions an individual cycles
through within a fiscal year --- is reported as
$\#\{r : (\text{id}_{i}, t, r) \in \mathcal{D}_{\mathrm{OTIS}}\}$
distributed across $i$.

\paragraph{Mortifying dimension.}
Goffman's mortification thesis predicts that institutional
processing produces co-occurring stigma markers. MRM tests this
by computing the joint Pearson $\chi^{2}$ on the alert-state
combinations $\{\mathrm{MH}, \mathrm{SR}, \mathrm{SW}\}$ and
their pairwise interactions, with Cram\'{e}r's $V$ as the effect
size. The disciplinary--medical overlap is the analogous
$\chi^{2}$ on the cross-tabulation
\[
  (\text{disciplinary alert}) \times (\text{medical-protective alert}).
\]
Statistically significant co-occurrence is taken as evidence of
a mortifying institutional process that cross-classifies the
individual along multiple stigma axes simultaneously.

\paragraph{Embedding dimension.}
Embedding is measured by the total-days survival distribution.
MRM fits both a log-normal and a Pareto distribution to the
aggregate-duration variable from \texttt{b02} and selects between
them by AIC:
\begin{equation}
  \mathrm{AIC}_{F}
  \;=\; -2\, \log\hat{L}_{F}
   \;+\; 2\, k_{F},
   \qquad F \in \{\text{log-normal}, \text{Pareto}\}.
  \label{eq:embedding-aic}
\end{equation}
A Pareto-best fit is consistent with a heavy-tailed embedding
process (a small number of individuals with extreme institutional
embedding); a log-normal-best fit is consistent with a more
gradual embedding gradient.

\paragraph{Intra-year transition matrix.}
The cyclic-region dimension is also captured by an intra-fiscal-
year first-order Markov transition matrix on the
five-region state space:
\begin{equation}
  P^{(t)}_{r \to r'}
  \;=\;
  \frac{\#\{(\mathrm{id}, t, r, r') :
              r \to r' \text{ within FY } t\}}
       {\#\{(\mathrm{id}, t, r) :
              r \text{ visited in FY } t\}}.
  \label{eq:transition-matrix}
\end{equation}
The diagonal entries $P^{(t)}_{r \to r}$ measure persistence;
off-diagonal mass measures cycling.

\paragraph{Path complexity.}
A within-cell Gini on the path-length distribution
(split by year $\times$ region) summarises how concentrated the
cycling is among a small set of high-churn individuals versus
distributed across the whole population.

\paragraph{Goffmanian-IRR via GLMM.}
The cycling effect on the count outcome is also estimated by a
generalised linear mixed model:
\begin{equation}
  \log \E[Y_{i}^{\mathrm{vm}} \mid X_{i}, u_{j(i)}]
  \;=\; \beta^{\T} X_{i} \;+\; u_{j(i)},
  \qquad u_{j} \sim \mathcal{N}(0, \sigma_{u}^{2}),
  \label{eq:churn-glmm}
\end{equation}
under both Poisson and negative-binomial families, with
$j(i)$ the institution containing individual $i$. The IRR for
treatment $T \in \{T^{\mathrm{ac}}, T^{\mathrm{any}}\}$
(\S\ref{sec:dual-rf}) is reported alongside its E-value.

\begin{remark}[Anonymisation caveat]
\label{rmk:churn-anon}
OTIS \texttt{UniqueIndividual\_ID} is anonymised as
\texttt{YYYY-XXXXX-AA}: the \texttt{YYYY} prefix locks each
identifier to one fiscal year, and identifiers are re-randomised
per dataset file even within the same year. The
churn analyses of \S\ref{sec:churn} are therefore strictly
\emph{within-year} estimands; I do not (and cannot) follow
individuals across fiscal years from the public OTIS aggregate.
A multi-year cohort design would require linkage to a
non-public per-person record, which is outside the scope of
MRM as currently released.
\end{remark}

\section{OTIS$\times$TPS overlay}
\label{sec:overlay}

MRM supports a cross-jurisdiction overlay between OTIS
provincial restrictive-confinement counts and TPS crime counts.
For a paired annual series of provincial Toronto-region
segregation counts $S_{t}$ and TPS Toronto-wide incident counts
$N^{(\mathrm{TPS})}_{t}(c)$ for category $c$, the year-over-year
correlation is
\begin{equation}
  \widehat{\rho}_{S, N(c)}
  \;=\;
   \frac{\sum_{t=2}^{T}
         (\Delta S_{t})(\Delta N^{(\mathrm{TPS})}_{t}(c))}
        {\sqrt{\sum_{t=2}^{T} (\Delta S_{t})^{2}\,
               \sum_{t=2}^{T} (\Delta N^{(\mathrm{TPS})}_{t}(c))^{2}}},
  \quad
  \Delta S_{t} \eqdef S_{t} - S_{t-1}.
  \label{eq:overlay-yoy}
\end{equation}
A composite-overlay statistic averages
$\widehat{\rho}_{S, N(c)}$ across the nine TPS categories,
weighting by category-specific CSI weights (\S\ref{sec:csi}). The
overlay is the framework's way of asking whether changes in
provincial restrictive confinement lead, lag, or co-vary with
changes in police-recorded crime --- an intentionally minimal
identification claim, but a useful descriptive backbone for
further causal work.

A per-region rollup
\begin{equation}
  \widehat{\rho}_{S, N(c)}^{(r)}
  \;=\;
   \frac{\sum_{t=2}^{T}
         (\Delta S^{(r)}_{t})(\Delta N^{(\mathrm{TPS})}_{t}(c))}
        {\sqrt{\sum_{t=2}^{T} (\Delta S^{(r)}_{t})^{2}\,
               \sum_{t=2}^{T} (\Delta N^{(\mathrm{TPS})}_{t}(c))^{2}}}
  \label{eq:overlay-region}
\end{equation}
disaggregates the YoY correlation by Ontario administrative
region $r \in \mathcal{R}$, surfacing regional heterogeneity that
the city-wide statistic averages over.

\section{Ontario Special Investigations Unit automation}
\label{sec:siu-ontario}

The Ontario Special Investigations Unit publishes
director's-report HTML pages keyed by an internal locator
identifier $\mathrm{drid}$. Each report has a paired news-release
locator $\mathrm{nrid}$. The two are not equal, and the canonical
join key is the SIU \emph{case number} of the form
$\mathrm{YY\text{-}XXX\text{-}NNN}$. The Ontario~SIU pipeline in
\texttt{morie.siu} ingests the corpus end-to-end:
\begin{equation}
  \texttt{scrape\_drid} \to \texttt{parse\_html}
  \to \texttt{normalize} \to \texttt{schema}
  \to \texttt{write\_csv} \to \texttt{analyze}.
  \label{eq:siu-pipeline}
\end{equation}
The schema $\Sigma_{\mathrm{SIU}}$ has 65 columns and is
case-number-keyed. The \texttt{analyze} module produces seven
RichResult-emitting summaries: \texttt{by\_police\_service},
\texttt{by\_year}, \texttt{case\_counts}, \texttt{demographics},
\texttt{mental\_health\_race\_indicators},
\texttt{decision\_timing}, \texttt{all\_analyses}. The pipeline
is conservative on the network side: $\mathrm{drid}$ ranges are
scraped politely with a concurrency limit and an HTTP cache, and
all analytical outputs are reproducible from the canonical CSV.

This pipeline is the framework author's contribution and was
designed in early 2026 from a brief draft titled
\emph{SIU\_initial\_draft.md}~\cite{Ruhela2026SIUDraft}.

\section{Federal aggregate companion: CCRSO and Doob T-539-20}
\label{sec:ccrso}

\subsection{Decoupling test}

For matched annual time series $C = (C_{t})_{t=1}^{T}$ of
crime rate and $I = (I_{t})_{t=1}^{T}$ of imprisonment rate, the
decoupling statistic is the Pearson correlation
\begin{equation}
  \widehat{\rho}_{C, I} \;=\;
  \frac{\sum_{t=1}^{T}
        (C_{t} - \bar{C})(I_{t} - \bar{I})}
       {\sqrt{\sum_{t=1}^{T}(C_{t}-\bar{C})^{2}
              \sum_{t=1}^{T}(I_{t}-\bar{I})^{2}}}.
  \label{eq:decoupling}
\end{equation}
A small or sign-mismatched $\widehat{\rho}_{C,I}$ supports the
Doob~\cite{Doob2020Affidavit} thesis that imprisonment rates
have decoupled from crime rates. MRM reports
$\widehat{\rho}_{C,I}$ alongside Pettitt's
non-parametric change-point statistic
\begin{equation}
  U_{T}(t) \;=\; \sum_{i=1}^{t} \sum_{j=t+1}^{T}
   \mathrm{sgn}\bigl(X_{i} - X_{j}\bigr),
  \quad
  K_{T} \;=\; \max_{1 \le t < T} \lvert U_{T}(t) \rvert,
  \label{eq:pettitt}
\end{equation}
applied to each of $C$ and $I$.

\subsection{CCRSO tables}

MRM hardcodes Tables 1, 2, and 3 from Doob's $T$-$539$-$20$
affidavit (Vol.\ 3, pp.\ 778--795, sourced from CCRSO 2018):
\begin{itemize}
  \item \textbf{Table 1} (5-year-average annual conditional
        releases by type): violent revocation rates of
        $0.14\%$ (day parole), $0.49\%$ (full parole), and
        $1.50\%$ (statutory release).
  \item \textbf{Table 2} (prisoner flow 2013/14--2017/18):
        average daily count $\approx 14{,}638$; statutory
        releases $\approx 5{,}127$/year, or about $427$/month.
  \item \textbf{Table 3} (age distribution 2018): age-group
        $50+$ comprises $47.3\%$ of the Canadian adult population
        but only $25.2\%$ of CSC custody and $16.9\%$ of CSC
        admissions.
\end{itemize}

\section{A worked example: \texttt{c02} institution-mix MRM}
\label{sec:worked-example}

To illustrate how MRM's pieces compose end-to-end, I walk
through the analysis of OTIS \texttt{c02}, an aggregate table of
restrictive-confinement counts indexed by institution $j$, fiscal
year $t$, region $r$, and gender $g$. I am interested in the
conditional incidence-rate ratio for treatment
$T = \one\{g = \text{Female}\}$ controlling for the
high-cardinality institution and region effects.

\paragraph{Step 1: aggregate model.} The starting point is the
log-linear fit~\eqref{eq:rf-aggregate-model} with stratum fixed
effects for institution. Because $J = 25$ institutions are too
many for a low-dimensional Poisson MLE and the institutional
denominator is variable across years, MRM instead fits the
marginal model~\eqref{eq:rf-gee-marginal} with cluster-robust
standard errors and an institution cluster.

\paragraph{Step 2: GEE grid.} MRM fits the four GEE
combinations of family $F \in \{\text{Pois}, \text{NB}\}$ and
working correlation $\mathbf{W} \in \{\text{Exch}, \text{Indep}\}$
in parallel, retaining the converged ones. Following
Theorem~\ref{thm:gee-priority}, the (NB, Exch) row is taken as the
canonical estimate; the others are reported for sensitivity. When
the NB-Exch fit fails to converge --- which occurs for some
\texttt{c}-series tables with extreme overdispersion --- MRM
falls back through (Pois, Exch), (NB, Indep), and (Pois, Indep)
in that order.

\paragraph{Step 3: report.} The output is a \texttt{RichResult}
with the fitted IRR, its $95\%$ Wald CI~\eqref{eq:irr-ci}, the
point estimate's $E$-value~\eqref{eq:evalue}, and a four-row
sensitivity table covering all converged GEE fits. Where
individual-level data are available (the \texttt{a}- and
\texttt{b}-series), the per-row estimators of
\S\ref{sec:mrm-perrow} are also reported.

\paragraph{Step 4: companion analyses.} MRM automatically pairs
the aggregate component with (i) the corresponding Doob
$\chi^{2}$ on the appropriate 2-way slice, (ii) the
cross-jurisdiction Mandela classifier for the same year and
region (via \texttt{c11}), (iii) the federal Sprott--Doob
comparison for the same fiscal period, and (iv) the
OTIS$\times$TPS YoY correlation for the same region. The
resulting \emph{master report} has six sections
(\S\S\ref{sec:mrm-aggregate}--\ref{sec:overlay}).

\section{Replication of Sprott--Doob--Iftene}
\label{sec:replication}

MRM's $\chi^{2}$ verifier rebuilds five published $\chi^{2}$
statistics from the transcribed cell counts of
Sprott--Doob--Iftene's three reports. The recomputed values
match the published values to within $0.01$:

\begin{table}[ht]
  \caption{$\chi^{2}$ verification (recomputed vs.\ published).}
  \label{tab:chi2-verify}
  \centering
  \begin{tabular}{lrrrr}
    \toprule
    Source                                      & Recomputed & df & Published & df \\
    \midrule
    SD 2021 Feb T11: Region $\times$ stay length & 201.01    & 16 & 201.00    & 16 \\
    SD 2021 Feb T15: Region $\times$ MH-flag    &  27.51    &  4 &  27.51    &  4 \\
    SD 2021 Feb T22: Region $\times$ Mandela    & 208.54    &  8 & 208.54    &  8 \\
    SDI 2021 May T5: Race $\times$ \#reviews    &   8.50    &  3 &   8.50    &  3 \\
    SDI 2021 May T10: Per-IEDM remain           &  26.12    & 11 &  26.12    & 11 \\
    \bottomrule
  \end{tabular}
\end{table}

The verification confirms both the accuracy of the framework's
transcription and the internal consistency of the published
statistics. I provide a generic
\texttt{verify\_chi2(observed)} helper so that the same diagnostic
can be applied to any future replication target.

\section{Empirical findings}
\label{sec:findings}

I report three headline findings.

\paragraph{F1 (Mandela cross-jurisdiction).}
The provincial Ontario Mandela$^{\ast}$ \emph{Torture}-classified
proportion (duration-only proxy applied to OTIS \texttt{c11}
Segregation individuals) rises monotonically from $12.5\%$ in
fiscal 2023 to $20.6\%$ in fiscal 2025, exceeding the federal
SIU $\mathrm{Torture}$-classified proportion of $9.9\%$
that Sprott and Doob~\cite{SprottDoob2021Torture} found across
$N = 1{,}960$ federal person-stays (November 2019 -- September
2020). The peak provincial-vs-federal gap is $+10.7$ percentage
points. This is robust to the conservative duration-only proxy
(\S\ref{sec:mandela}).

\paragraph{F2 (IEDM oversight is structurally limited).}
I confirm the headline IEDM-process findings of
Sprott--Doob--Iftene~\cite{SprottDoobIftene2021IEDM} on $N = 265$
SIU person-stays receiving $\geq 1$ IEDM review under
CCRA~s.~$37.8$. Of $380$ rendered decisions, $58.9\%$ kept the
inmate in the SIU and $30.3\%$ were pre-empted by CSC moving the
prisoner before the IEDM ruled; only $8.7\%$ ordered the inmate's
removal. Twelve IEDMs varied from $37.5\%$ to $85.7\%$ in their
``should remain'' rate ($\chi^{2} = 26.12$, df $= 11$, $p < .01$),
indicating substantial inter-IEDM heterogeneity that is not
explained by the cases reviewed. $105$ stays of $\geq 76$ days,
including $49$ of the $137$ stays $\geq 120$ days, had no IEDM
record at all in the dataset.

\paragraph{F3 (Sprott--Doob--Iftene replication).}
All five published $\chi^{2}$ statistics from the three
Sprott--Doob--Iftene reports reproduce to within $0.01$ of the
published values from the transcribed cell counts
(Table~\ref{tab:chi2-verify}). The Pacific torture rate is
$22.6\times$ the Ontario federal rate ($39.1$ vs.\ $1.73$ per
$1{,}000$ regional prisoners; SD~2021~Feb~Table~23).

\section{Limitations and ethical considerations}
\label{sec:limitations}

\paragraph{Aggregation and identifiability.} OTIS is released as
an aggregate file. Per-individual identifiers are not available
beyond a year-bounded \texttt{unique\_individual\_id}, which means
the \texttt{a}- and \texttt{b}-series RFs in
\S\ref{sec:mrm-perrow} treat row-aggregations as analytic units
rather than person-time records. This is conservative: variances
are likely smaller than they would be in a true per-person
panel, but the point estimates are not biased provided the
within-aggregation correlation structure is well-modelled by the
GEE grid of \S\ref{sec:gee-grid}.

\paragraph{Mandela duration-only proxy.} As I am explicit in
\S\ref{sec:mandela}, the OTIS Mandela classification is a
duration-only proxy. The federal CSC SIU classification
operationalised by Sprott and Doob~\cite{SprottDoob2021Torture}
also requires the inmate to have spent $\le 2$ hours out-of-cell
on average and to have missed the four-hour requirement on
$100\%$ of days; OTIS does not expose hours-out-of-cell, so the
provincial proportion classified as
\emph{Torture}$^{\ast}$ may overstate the proportion that would
be classified \emph{Torture} under the full operationalization.
I therefore report two views (Segregation and Restrictive
Confinement) and treat the cross-jurisdiction gap as an upper
bound rather than a point estimate.

\paragraph{Two-SIU disambiguation.} As noted in
Remark~\ref{rmk:two-siu}, the federal Structured Intervention
Unit and the Ontario Special Investigations Unit share the
acronym SIU but refer to entirely distinct agencies. MRM's
MORIE instantiation deliberately uses \texttt{morie.siuiap}
for the former and \texttt{morie.siu} for the latter; readers
and downstream code should never conflate them.

\paragraph{Data privacy.} The Ontario SIU scraper of
\S\ref{sec:siu-ontario} only ingests publicly published
director's-report HTML and the paired news-release HTML. No
personally identifying information beyond what is published is
collected, and all output CSVs are case-number-keyed (the SIU's
own canonical identifier).

\paragraph{Non-causal identification claims.} The OTIS$\times$TPS
overlay (\S\ref{sec:overlay}) and the Lotka--Volterra dynamics
(\S\ref{sec:tps-physics}) are descriptive and do not, on their
own, license causal-identification claims. The per-row component
of \S\ref{sec:mrm-perrow} is the only set of estimators in MRM
that targets a formal causal estimand, and only under the
standard SUTVA, consistency, positivity, and unconfoundedness
assumptions.

\section{Reproducibility}
\label{sec:reproducibility}

MRM is implemented in the open-source MORIE package as a
set of MRM modules, distributed via Python and R. At the
time of this paper:
\begin{itemize}
  \item \textbf{Tests.} The Python test suite contains $\approx
        14{,}836$ baseline tests plus $55$ tests of the federal
        / provincial Mandela cross-jurisdiction modules, all
        passing.
  \item \textbf{Documentation.} A Sphinx documentation site
        builds without warnings under the strict-mode
        (\verb|-W|) build, including pages for each component
        described in this paper.
  \item \textbf{Datasets.} The 41 built-in datasets include the
        OTIS aggregate file, the Ontario SIU canonical CSV, the
        TPS open-data categories, and the CCRSO 2018 table
        constants.
\end{itemize}

The package is released under the GNU General Public License,
version 2 only (\textbf{GPL-2.0-only};
\url{https://www.gnu.org/licenses/old-licenses/gpl-2.0.html}), with
the Ontario~SIU ingestion utilities additionally subject to the
Ontario open-government licence on the source corpus. The paper
text and accompanying figures are released under \textbf{CC-BY-4.0}
(\url{https://creativecommons.org/licenses/by/4.0/}), the licence
the Zenodo record at \url{https://doi.org/10.5281/zenodo.20096075}
carries.

\section*{Acknowledgements}

The author thanks Prof.\ Beatrice Jauregui (University of Toronto)
for the network of mentorship that made this line of work
possible, and Prof.\ Ayobami Laniyonu (University of Toronto) for
valuable lessons over the broader period of this work that
informed the author's thinking on policing and corrections
research. The MRM methodology lineage acknowledges the federal
context provided by Anthony N.\ Doob's affidavit and the four
Sprott--Doob--Iftene independent academic reports
(\S\ref{sec:mrm}). All implementation, all framework design, and
all empirical findings reported in this paper are the work of the
framework author.

\bibliographystyle{plainnat}
\begin{thebibliography}{99}

\bibitem{Bettencourt2007}
Bettencourt, L.\ M.\ A., Lobo, J., Helbing, D., K\"uhnert, C., \&
West, G.\ B.\ (2007).
\emph{Growth, innovation, scaling, and the pace of life in
cities.}
Proc.\ Natl.\ Acad.\ Sci.\ USA, 104(17), 7301--7306.

\bibitem{Daley2003PointProcess}
Daley, D.\ J., \& Vere-Jones, D.\ (2003).
\emph{An Introduction to the Theory of Point Processes,
Volume I: Elementary Theory and Methods} (2nd ed.).
Springer.

\bibitem{Chernozhukov2018DoubleML}
Chernozhukov, V., Chetverikov, D., Demirer, M., Duflo, E.,
Hansen, C., Newey, W., \& Robins, J.\ (2018).
\emph{Double/debiased machine learning for treatment and
structural parameters.}
The Econometrics Journal, 21(1), C1--C68.

\bibitem{ClausetShalizi2009Powerlaw}
Clauset, A., Shalizi, C.\ R., \& Newman, M.\ E.\ J.\ (2009).
\emph{Power-law distributions in empirical data.}
SIAM Review, 51(4), 661--703.

\bibitem{Goffman1961Asylums}
Goffman, E.\ (1961).
\emph{Asylums: Essays on the Social Situation of Mental Patients
and Other Inmates.}
Anchor Books.

\bibitem{Doob2020Affidavit}
Doob, A.\ N.\ (2020).
\emph{Affidavit of Anthony N. Doob} (Federal Court of Canada,
File T-539-20, Application Record Vol.\ 3 of 5, pp.\ 778--795).
Canadian Civil Liberties Association et al.\ v.\ Attorney General
of Canada.

\bibitem{Ruhela2026MRM}
Ruhela, V.\ S.\ (2026).
\emph{The MRM Framework: A Multi-Source Statistical Foundation
for Canadian Carceral, Police, and Oversight Data, Implemented
as MRM modules in MORIE.}
Zenodo. \url{https://doi.org/10.5281/zenodo.20096075} (this paper).

\bibitem{Ruhela2026MORIE}
Ruhela, V.\ S.\ (2026).
\emph{MORIE: A Multi-Domain Scientific Computing Toolkit for
Observational Inference, with Sociolegal, Signal-Processing,
Cryptographic, and Spatial-Statistics Modules.}
Zenodo. \url{https://doi.org/10.5281/zenodo.20096350}.

\bibitem{Ruhela2026Software}
Ruhela, V.\ S.\ (2026).
\emph{MORIE Toolkit: Methods for Observational Inference and
Robust Analysis of Interventions in Scientific Experimentation}
[Software, v0.1.2].
Zenodo. \url{https://doi.org/10.5281/zenodo.20111233}.
PyPI: \url{https://pypi.org/project/morie/}.
r-universe: \url{https://hadesllm.r-universe.dev/morie}.
Documentation: \url{https://hadesllm.github.io/morie/}.

\bibitem{Ruhela2026Hawkes}
Ruhela, V.\ S.\ (2026).
\emph{Criminological Hawkes Process via MORIE: Markovian and
Non-Markovian Self-Exciting Point Processes for Toronto Crime.}
Zenodo. \url{https://doi.org/10.5281/zenodo.20102198}.

\bibitem{Ruhela2026SIUDraft}
Ruhela, V.\ S.\ (2026).
\emph{SIU initial draft (Ontario Special Investigations Unit
automation).}
Internal manuscript, May 2026.

\bibitem{ShortBrantingham2008}
Short, M.\ B., D'Orsogna, M.\ R., Pasour, V.\ B., Tita, G.\ E.,
Brantingham, P.\ J., Bertozzi, A.\ L., \& Chayes, L.\ B.\ (2008).
\emph{A statistical model of criminal behavior.}
Mathematical Models and Methods in Applied Sciences, 18(suppl.),
1249--1267.

\bibitem{SprottDoob2020Operation}
Sprott, J.\ B., \& Doob, A.\ N.\ (2020, October).
\emph{Understanding the Operation of Correctional Service
Canada's Structured Intervention Units: Some Preliminary
Findings.}
Centre for Criminology \& Sociolegal Studies, University of
Toronto.

\bibitem{SprottDoob2020Covid}
Sprott, J.\ B., \& Doob, A.\ N.\ (2020, November).
\emph{Is there Clear Evidence that COVID-19 Was the Cause of
Problems with the Operation of CSC's Structured Intervention
Units?}
Centre for Criminology \& Sociolegal Studies, University of
Toronto.

\bibitem{SprottDoob2021Torture}
Sprott, J.\ B., \& Doob, A.\ N.\ (2021, February).
\emph{Solitary Confinement, Torture, and Canada's Structured
Intervention Units.}
Centre for Criminology \& Sociolegal Studies, University of
Toronto.

\bibitem{SprottDoobIftene2021IEDM}
Sprott, J.\ B., Doob, A.\ N., \& Iftene, A.\ (2021, May).
\emph{Do Independent External Decision Makers Ensure that ``An
Inmate's Confinement in a Structured Intervention Unit Is to End
as Soon as Possible''?}
Schulich School of Law, Dalhousie University.

\bibitem{UN2015Mandela}
United Nations General Assembly (2015).
\emph{United Nations Standard Minimum Rules for the Treatment of
Prisoners (the Nelson Mandela Rules).}
A/Res/70/175.

\bibitem{vanderLaan2007SuperLearner}
van der Laan, M.\ J., Polley, E.\ C., \& Hubbard, A.\ E.\ (2007).
\emph{Super Learner.}
Statistical Applications in Genetics and Molecular Biology, 6(1).

\bibitem{VanderWeele2017Evalue}
VanderWeele, T.\ J., \& Ding, P.\ (2017).
\emph{Sensitivity analysis in observational research:
Introducing the E-value.}
Annals of Internal Medicine, 167(4), 268--274.

\bibitem{Wallace2009CSI}
Wallace, M., Turner, J., Matarazzo, A., \& Babyak, C.\ (2009).
\emph{Measuring crime in Canada: Introducing the Crime Severity
Index and improvements to the Uniform Crime Reporting Survey.}
Statistics Canada Catalogue no.\ 85-004-X.

\bibitem{ZorroMedina2023JMP}
Zorro Medina, {\'A}.\ (2023).
\emph{The Effect of Anti-Gang Laws on Crime and Social Control.}
Job Market Paper, University of Chicago.
\url{https://azorromedina.com/wp-content/uploads/2023/08/JMP_ZorroMedina_28_08_23H.pdf}.

\bibitem{Athey2022Staggered}
Athey, S., \& Imbens, G.\ W.\ (2022).
\emph{Design-based analysis in difference-in-differences
settings with staggered adoption.}
Journal of Econometrics, 226(1), 62--79.

\bibitem{CallawaySantanna2021}
Callaway, B., \& Sant'Anna, P.\ H.\ C.\ (2021).
\emph{Difference-in-differences with multiple time periods.}
Journal of Econometrics, 225(2), 200--230.

\bibitem{GoodmanBacon2021DiD}
Goodman-Bacon, A.\ (2021).
\emph{Difference-in-differences with variation in treatment
timing.}
Journal of Econometrics, 225(2), 254--277.

\end{thebibliography}

\end{document}
